%%%%%%%%%%%%%%%%%%%%%%%%%%%%%%%%%%%%%%%%%%%%%%%%%%%%%%%%%%%%%%%%%%%%%
%%                                                                 %%
%% Please do not use \input{...} to include other tex files.       %%
%% Submit your LaTeX manuscript as one .tex document.              %%
%%                                                                 %%
%% All additional figures and files should be attached             %%
%% separately and not embedded in the \TeX\ document itself.       %%
%%                                                                 %%
%%%%%%%%%%%%%%%%%%%%%%%%%%%%%%%%%%%%%%%%%%%%%%%%%%%%%%%%%%%%%%%%%%%%%
\documentclass[12pt]{article} 

\usepackage{authblk} \usepackage{amsmath, amsfonts,
	mathpazo} \usepackage{geometry}

\usepackage{cite} \usepackage{graphicx} \usepackage{amsmath}
\usepackage{algorithm,algorithmicx} % \usepackage{undefinedalgorithmic}
\usepackage{algpseudocode} \usepackage{graphics} \usepackage{epsfig}

\usepackage{comment} \usepackage{subfigure}
\usepackage{url}
\usepackage{tikz} \usetikzlibrary{positioning} %为了实现相对位置的设定
\usepackage{xcolor} %为了实现不同的颜色\usepackage{pgfplots}
\usepackage{pstricks}

\definecolor{mycolor}{rgb}{0.1,0.5,1}

\newtheorem{remark}{Remark} \newtheorem{example}{Example}
\newtheorem{approach}{Approach}
\usepackage{booktabs}
\newcommand\keywords[1]{\textbf{Keywords}: #1}

\geometry{ a4paper, margin=0.5in }

\usepackage{multicol}
\numberwithin{figure}{section}
\numberwithin{table}{section}



%% as per the requirement new theorem styles can be included as shown below
\newtheorem{theorem}{Theorem}%  meant for continuous numbers
%%\newtheorem{theorem}{Theorem}[section]% meant for sectionwise numbers
%% optional argument [theorem] produces theorem numbering sequence instead of independent numbers for Proposition
\newtheorem{proposition}[theorem]{Proposition}% 
%%\newtheorem{proposition}{Proposition}% to get separate numbers for theorem and proposition etc.
%\newtheorem{example}{Example}%
%\newtheorem{remark}{Remark}%
%\newtheorem{definition}{Definition}%
\usepackage{listings}
\usepackage{authblk}

\raggedbottom
%%\unnumbered% uncomment this for unnumbered level heads

\title{An orthogonalization-free implementation of the LOBPCG method in solving Kohn-Sham equation} 

\author[1]{Chengyu Liu} \author[1,2,3]{Guanghui Hu}

\affil[1]{Department of Mathematics, Faculty of Science and
	Technology, University of Macau, Macao SAR, China.}
\affil[2]{Zhuhai UM Science \& Technology Research Institute, Zhuhai,
	Guangdong, China.}
\affil[3]{Guangdong-Hong Kong-Macao Joint Laboratory for Data-Driven
	Fluid Mechanics and Engineering Applications, University of Macau,
	Macao SAR, China.}
%%==================================%%
%% sample for unstructured abstract %%
%%==================================%%


\begin{document}

\maketitle
\abstract{In the classic implementation of the LOBPCG method, orthogonalization and the R-R (Rayleigh-Ritz) procedure cost nonignorable CPU time. Especially this consumption could be very expensive to deal with situations with large block sizes. In this paper, we propose an orthogonalization-free framework of implementing the LOBPCG method for SCF (self-consistent field) iterations in solving the Kohn-Sham equation. In this framework, orthogonalization is avoided in calculations, which can decrease the computational complexity. And the R-R procedure is implemented parallelly through OpenMP, which can further reduce computational time. During numerical experiments, an effective preconditioning strategy is designed, which can accelerate the LOBPCG method remarkably. Consequently, the efficiency of the LOBPCG method can be significantly improved. Based on this, the SCF iteration can solve the Kohn-Sham equation efficiently. A series of numerical experiments are inducted to demonstrate the effectiveness of our implementation, in which significant improvements in computational time can be observed.}

\keywords{LOBPCG method; orthogonalization-free; Kohn-Sham equation; preconditioning strategy.}
\section{Introduction}\label{sec1}
The LOBPCG method \cite{knyazev2001toward} has been widely used for solving generalized eigenvalue problems in comprehensive domains due to its efficiency \cite{yamada200516,yamada2006high,yang2009kssolv,bao2012h,mcqueen2016megaman,dziekonski2017gpu,shao2018accelerating}. Especially in electronic structure calculations, LOBPCG method plays an important role in the SCF iteration for solving Kohn-Sham equation \cite{lin2012adaptive,bao2012h,bao2013numerical,lin2019numerical,knyazev2014numerical}.

However, the LOBPCG method still encounters some challenges when dealing with large-scale problems since the use of orthogonalization and Rayleigh-Ritz procedure will cost much CPU time. As mentioned by Knyazev in \cite{knyazev2017recent}, to further improve the efficiency of the LOBPCG method, more efforts are needed to decrease the use of orthogonalizations and Rayleigh-Ritz, while preserving fast convergence, maintaining reasonable stability, and keeping the algorithm simple to implement. This limitation comes from the time complexity of the RR procedure (at least $\mathcal{O}(Np)$) and the orthogonalization (at least $\mathcal{O}(Np^2)$), where $N$ denotes the dimensionality and $p$ is the number of computed eigenvectors. Thus, the computational time for solving large-scale problems can be quite large.

To overcome such a limitation, the Rayleigh-Ritz procedure can be efficiently implemented through some parallel libraries due to the block structure in the algorithm \cite{bottin2008large,duersch2018robust}. In this way, the time complexity can be reduced to around $\mathcal{O}(N)$ ideally. Thus, the calculation time will be reduced considerably. But the parallelization of orthogonalization is still notoriously difficult by now. So it is impossible to gain acceleration of orthogonalization through the parallelization. To deal with this problem, researchers develop some other techniques to improve the efficiency of orthogonalization such as \cite{dziekonski2017gpu,kressner2023mixed}. However, the time complexity of orthogonalization is still $\mathcal{O}(Np^2)$. Consequently, the time consumption of orthogonalization is still unignorable during calculations of large-scale problems.

As for ensuring fast convergence and reasonable stability of the LOBPCG method, it can be achieved by some suitable preconditioners. There have been many preconditioners proposed for the LOBPCG method, such as incomplete Cholesky factorization, Davidson, filter preconditioners and so on \cite{knyazev2001toward,saad1996solution,arbenz2012lecture,seitsonen1995real,GAPP}. Typically, Poisson preconditioner using the inverse of the kinetic operator as a specific preconditioner for the Kohn-Sham equation is also a possible choice \cite{EarlyPoissonPre,Octopus}. But the time consuming is a problem of this preconditioner. 

In this paper, based on the adaptive finite element method framework, we design an efficient implementation of LOBPCG methods for SCF (self-consistent field) iteration in KS-DFT (Kohn-Sham density functional theory). In our implementation, the orthogonalization in the LOBPCG method is avoided, significantly improving the efficiency without affecting the accuracy of the numerical results. Furthermore, Rayleigh-Ritz procedure is implemented in parallelization through OpenMP \cite{dagum1998openmp}. Consequently, it can be expected that the computational complexity of the whole calculation is close to $\mathcal{O}(N)$. Some related numerical experiments are conducted to confirm the feasibility of our implementation. The CPU time is saved significantly while numerical results' accuracy is not influenced much without orthogonalization in LOBPCG. Besides that, we give some potential explanations of why such an orthogonalization-free implementation of the LOBPCG method can work. 

In addition, a mixing preconditioning strategy is used in our implementation to ensure the fast convergence of the LOBPCG method. In the case of some complex molecules, such as benzene, it isn't easy to obtain results to the specified accuracy with the Davidson preconditioner. Then we turn to use Poisson preconditioners, which can lead to an accurate enough result. However, the CPU time for calculating the inverse of the kinetic operator is unaffordable. To deal with this, we combine Davidson and Poisson preconditioners for the LOBPCG method, which lets us get accurate results with acceptable time consumption. 

Moreover, based on AFEPack library \cite{li2005multi}, we implement a series of calculations for solving the KS equation of different molecules. With the efficient implementation of the LOBPCG method, the SCF iteration performs better convergence. As a result, the efficiency of the whole calculation is improved remarkably. In this way, accurate enough ground state energy of specific molecules can be calculated with acceptable time consumption. 

The paper is organized as follows. In Section 2, the classic framework of LOBPCG algorithm in SCF iteration for solving KS equation is reviewed. Then the efficient implementation of LOBPCG method for SCF iterations in KS-DFT is introduced in detail. The improvement in orthogonalization and Rayleigh-Ritz procedure is discussed in Section 3. Then, in Section 4, numerical examples are provided to demonstrate the effectiveness and acceleration of this implementation for solving the Kohn-Sham equation. Finally, some concluding remarks and our future work are given in the last section.

\section{Review of LOBPCG}
Consider the Kohn-Sham equation of an $N$ particle system as follows:
\begin{equation}
\begin{aligned}
	\label{KS}
	H[n]\varphi_j(\boldsymbol{r})  = \varepsilon_j \varphi_j(\boldsymbol{r}),
\end{aligned}
\end{equation}
where the $\varepsilon_j$ are the Kohn-Sham eigenvalues and the $\varphi_j$ are named as Kohn-Sham orbitals or wavefunctions. By summing the squares of the lowest $N$ occupied orbitals, the electron density $n(\boldsymbol{r})$ is given by
\begin{equation}
	n(\boldsymbol{r}) = \sum_{j=1}^N \vert\varphi_j(\boldsymbol{r})\vert^2.
\end{equation}
With the electron density, the Kohn-Sham Hamiltonian can be expressed as
\begin{equation}
\begin{aligned}
	H[n]&=\left(-\frac{\nabla^2}{2}+v_s[n](\boldsymbol{r})\right),\\
	\centering v_s[n](\boldsymbol{r})& = v(\boldsymbol{r}) +v_{H}[n](\boldsymbol{r})+v_{xc}[n](\boldsymbol{r}),\\
	v_{H}[n](\boldsymbol{r})&=\int \frac{n(\boldsymbol{r})}{\vert \boldsymbol{r}- \boldsymbol{r}' \vert}d\boldsymbol{r},\\
	v_{xc}[n](\boldsymbol{r})&=\frac{\delta E_{xc}[n]}{\delta n(\boldsymbol{r})},
\end{aligned}
\end{equation}
where $v$ represents the external potential, $v_H$ is the Hartree potential, and $v_{xc}$ is the xc potential.

Due to the nonlinearity of (\ref{KS}), it is usually solved by iterative methods.  The SCF iterations method is the most popular one. In each iteration, the eigenvalue problem (\ref{KS}) has to be solved to update Kohn-Sham orbitals, which can be used to calculate the electron density for the following iteration. In fact, the Kohn-Sham equation can be considered as a generalized eigenvalue problem as follows

\begin{equation}
	\label{GeneralEig}
	A u = \lambda Bu(A,B\in\mathbb{R}^{n\times n}, u\in \mathbb{R}^n),
\end{equation}
where $A$ is Hermitian and $B$ is Hermitian positive definite. To solve (\ref{GeneralEig}), the LOBPCG method is one of the most popular choice. 

The basic idea of the LOBPCG algorithm developed by Knyazev can be summarized as follows. Firstly, eigenvalues of the pencil $(A,B)$ are arranged in an increasing order, which are denoted by $\lambda_1\leq\lambda_2\leq\dots\leq\lambda_n$. The corresponding eigenvectors are defined as $x_1,x_2,\dots,x_n$. If the first $m\leq n$ eigenpairs are required, the eigenvalue problem can be transformed as
\begin{equation}
	AX=BX\Lambda, \text{where } X=[x_1,x_2,\dots,x_m], \Lambda=diag([\lambda_1,\lambda_2,\dots,\lambda_m])
\end{equation}
Then the (i+1)-th approximation of $X$ can be expressed by
\begin{equation}
	X_{i+1} = [X_{i},W_{i},P_{i}]\begin{bmatrix}C^1_{i+1}\\C^2_{i+1}\\C^3_{i+1}\end{bmatrix}=S_i C_{i+1},
\end{equation}
where $W_i=T^{-1}(AX_i-X_i\Theta_i)$ with the preconditioner $T$ and $\Theta_i = X_i^* AX_i$. $P_i$ is an update direction obtained from previous searches recursively, which is given by
\begin{equation}
	P^{i+1}=[W^i,P^i]\begin{bmatrix}C_{i+1}^2\\C_{i+1}^3\end{bmatrix}.
\end{equation}

Coefficient matrices $C$ are determined at each iteration step, by solving the trace minimization problem
\begin{equation}
	\label{mintrace}
	\min_{X_{i+1}^* BX_{i+1}=I} tr(X_{i+1}^* AX_{i+1}) =  \min_{C_{i+1}^{*}S_{i}^{*}BS_{i}C_{i+1}=I} tr(C_{i+1}^* S_i^* AS_i C_{i+1}),
\end{equation}
where $tr(\cdot)$ represents the trace of a matrix. By using teh Rayleigh-Ritz procedure, the solution $C_{i+1}$ of (\ref{mintrace}) can be obtained from the eigenvectors belonging to the $m$ smallest eigenvalues of the generalized eigenvalue problem 
\begin{equation}
	(S_{i}^* A S_{i}) C_{i+1} = (S^{*}_{i}BS_{i}) C_{i+1}\Theta_{i+1},
\end{equation}
where $\Theta_{i+1}$ is a diagonal matrix, whose diagonal entries are approximations to the desirable eigenvalues. Based on these steps, the LOBPCG algorithm can be summarized as Algorithm \ref{LOBPCGAlgorithm}.

%\begin{algorithm}
%	\label{LOBPCGAlgorithm}
%	\caption{LOBPCG}
%	Require: Initial approximate eigenvectors $\boldsymbol{\tilde{X}}\in\mathbb{R}^{n\times m}$, the number of desired eigenpairs $k\leq m$;\\
%	\item[1.] Select an initial guess $\tilde{X}\in\mathbb{R}^{n\times m}$,
%	\item[2.] $\boldsymbol{X}_0 = \tilde{\boldsymbol{X}} \boldsymbol{Y}$ where $(\boldsymbol{C},\boldsymbol{\Theta}_0) = RR(\tilde{\boldsymbol{X}},m)$,
%	%2. $\boldsymbol{\Theta}=X_0^*AX_0$.\\
%	\item[3.] $\boldsymbol{P}=[]$,
%	\item[4.] For $k=0,1,2,\dots$ do
%	\item[5.]\quad $\boldsymbol{R}=\boldsymbol{AX}_k - \boldsymbol{BX}_k \boldsymbol{\Theta}_k$,
%	\item[6.]\quad Determine the number of converged eigenpairs $n_c$,
%	\item[7.]\quad If $n_c\geq p$, then 
%	\item[8.]\quad\quad return $\boldsymbol{X}=\boldsymbol{X}_k(:,1:p)$ and $\boldsymbol{\Theta} = \boldsymbol{\Theta}_k(1:p,1:p)$;
%	\item[9.]\quad Solve the preconditioned linear system $\boldsymbol{TH=R}$.
%	\item[10.]\quad Let $\boldsymbol{S}=[\boldsymbol{X}_k,\boldsymbol{H},\boldsymbol{P}]$ and compute $(\boldsymbol{C},\boldsymbol{D}_{k+1}) = RR(\boldsymbol{S},m)$, i.e., solving the eigenvalue problem $\boldsymbol{A}_p \boldsymbol{C} = \boldsymbol{CD}$, $\boldsymbol{A}_p = \boldsymbol{S^*AS}$
%	\item[11.]\quad Orthogonalize  $\boldsymbol{C}(1:m,m+1:n)^*$ to obtain an orthogonal matrix $V$,
%	\item[12.]\quad $\boldsymbol{X}_k=\boldsymbol{SC}(1:m),\boldsymbol{\Theta}_k=\boldsymbol{D}(1:m,1:m),\boldsymbol{P}=SC(:,m+1:2m)V$.
%	\item[13.] Endfor 
%\end{algorithm}


\begin{algorithm}[htbp]
	\caption{The LOBPCG algorithm}
	\textbf{Input:} \\
	$X_0\in\mathbb{R}^{n\times m}$ is a matrix of initial approximate eigenvectors,\\ 
	$k$, the number of desired eigenpairs, $k\leq m$,\\
	maxiter, the maximal number of iterations.\\
	\textbf{Output:} \\
	\quad $\Theta$, the diagonal matrix of computed smallest eigenvalues, \\
	\quad $X$, the matrix containing the corresponding computed eigenvectors satisfying $AX = X\Theta$.
	\begin{algorithmic}[1]
		%\State $X_0 = \tilde{X} C_1$ where $(C_1,\Theta_1) = RR(\tilde{X},m)$,
		\State $P_0=[]$, $\Theta_0 = (X_0^* AX_0)/(X_0^* B X_0)$;
		\For{$k=0,1,2,\dots$, maxiter} 
		\State $R_k=AX_k - BX_k \Theta_k$;
		\State Determine the number of converged eigenpairs $n_c$;
		\If{$n_c > k$}
		\State Return $X=X_k$, $\Theta = \Theta_k$;
		\EndIf
		\State Solve the preconditioned linear system $T_k W_k=R_k$;
		\State $S_{k}=[X_k,W_k,P_k]$;
		\State Compute $(C_{k+1},\Theta_{k+1}) = RR(S_k,m)$;
		\State Orthogonalize  $C_{k+1}(1:m,m+1:n)^*$ to obtain an orthogonal matrix $V$;
		\State $X_{k+1}=S_k C_{k+1}(1:m),\Theta_{k+1}=\Theta_{k+1}(1:m,1:m),P_{k+1}=S_kC_{k+1}(:,m+1:2m)V$.
		\EndFor
	\end{algorithmic}
	\label{LOBPCGAlgorithm}
\end{algorithm}

With the LOBPCG method, we can obtain a numerical solution of (\ref{KS}) in each SCF iteration. Through SCF iteration, an approximate density can be obtained, with which the ground state energy can be calculated. However, the numerical result may be inaccurate due to the finite element space limit. To deal with this, we refine the mesh grid to build a more accurate space through AFEM (adaptive finite element method). The details of the mesh adaptation can be found in \cite{bao2012h,bao2013numerical}.The whole framework of the AFEM-based SCF iterations in solving Kohn-Sham equations is summarized in the Algorithm \ref{SCFinKSDFT}. It can be seen that the Kohn-Sham equation is solved by the LOBPCG method in each SCF iteration. Therefore, if the LOBPCG method is effectively implemented, the efficiency of the entire SCF iteration can be expected to increase accordingly. Based on this, we focusing on accelerating the calculation of the LOBPCG method.

\begin{algorithm}[htbp]
	\caption{AFEM based SCF iteration for solving KS equations}
	\textbf{Input:}\\ 
	$\{\varphi_j^0\}_{j=1}^{N},\dots,N$, the initial computed KS orbitals,\\
	maxadapt, the maximal number of mesh adaptation,\\
	maxiter, the maximal number of iterations,\\
	scftol, the tolerance of electronic density,\\
	Etol, the tolerance of computed ground state energy.\\
	\textbf{Output:}\\
	$\{\varphi_j\}_{j=1}^{N}$, the approximate KS orbitals,\\
	$E$, the approximate ground state energy.\\
	\begin{algorithmic}[1]
		%\State $X_0 = \tilde{X} C_1$ where $(C_1,\Theta_1) = RR(\tilde{X},m)$,
		\State $n^0=\sum_{j=1}^N \vert\varphi^0_j(\boldsymbol{r})\vert^2$;
		\State Update ground state energy $E^0$;
		\For{$i=0,1,2,\dots$, maxadapt}
		\For{$k=0,1,2,\dots$, maxiter} 
		\State Compute the Hartree $v_H^k$ based on $n^k$;
		\State Update Hamiltonian $H^k$;
		\State Solving the Kohn-Sham equation $H^k \varphi_j^{k+1}=\varepsilon_j^{k} \varphi_j^{k+1}$ through LOBPCG;
		\State  $\tilde{n}^{k} = \sum_{j=1}^N \vert\varphi^{k+1}_j(\boldsymbol{r})\vert^2$;
		\State Update the electron density through linear mixing: $n^{k+1}=n^{k}+\alpha(n^k-\tilde{n}^{k})$;
		\If{$\vert \frac{n^{k}-n^{k+1}}{n^{k}}\vert$ $<$ scftol}
		\State Break;
		\EndIf
		\EndFor
		\State Update ground state energy $E^k$;
		\If{$\vert \frac{E^{k}-E^{k-1}}{E^{k-1}}\vert<$  Etol}
		\State Break;
		\Else 
		\State Adapt mesh grid based on the computed Kohn-Sham orbitals; 
		\EndIf
		\EndFor
	\end{algorithmic}
	\label{SCFinKSDFT}
\end{algorithm}


\section{The orthogonalization-free implementation of LOBPCG}
To improve the efficiency of the LOBPCG method, we follow by the comment given by Knyazev, focusing on accelerating the orthogonalization and Rayleigh-Ritz procedure. In our numerical experiments of electronic structure calculations, the orthogonalization is avoided, and the Rayleigh-Ritz procedure is accelerated through parallelization. %Besides that, we design a mixing preconditioning strategy to improve the efficiency further. In the following, we discuss the feasibility and effectiveness of our implementation based on a series of numerical results.
\subsection{Orthogonalization}
Generally, orthogonalization is needed for ensuring the numerical stability of the LOBPCG \cite{hetmaniuk2006basis,duersch2018robust,kressner2023mixed}. However, orthogonalization is also a barrier to efficiency because its time complexity is $\mathcal{O}(Np^2)$. Furthermore, orthogonalization is extremely difficult to parallelize, which means it is impossible to reduce the time complexity by using more threads in calculations. Thus, the corresponding computational time will be nonignorable in large-scale cases. 

To gain more improvements in efficiency, in our implementation, orthogonalization is only applied on the initial eigenvectors that are generated randomly to ensure their $B$-orthogonality in the beginning of the whole iteration, i.e. at the step 1 of Algorithm \ref{SCFinKSDFT}. Then, in the following calculation, there is no more orthogonalization operation. %Without orthogonalization, the convergence is not influenced dramatically in our calculations, which can be observed from the following numerical experiments. 

There are several reasons leading to the feasibility  of our orthogonalization-free implementation, which are listed as follows. Firstly, the LOBPCG method is generally stable and robust. In the original implementation of LOBPCG, the orthogonalization is performed for very ill-conditioned
problems or when a very high accuracy is required \cite{knyazev2001toward}. To avoid these situations, we use shift-and-invert preconditioners in calculations. The computed eigenvalues are used as the approximation of exact eigenvalues to generate preconditioners during iterations. Through the shift-and-invert technique in LOBPCG, the computed eigenpairs gets close to the exact eigenvectors (wavefunctions) faster. 

In addition, the adaptive finite element method can capture the feature of the wavefunction more precisely. As a result, the discretized space can be more similar to the exact space, on which the numerical solution can be expected to be get close to the exact wavefunction. Physically, the exact orbitals (wavefunctions) are naturally orthogonal. When the calculated eigenvectors are sufficiently accurate, there will naturally be orthogonality between them. Combining the above reasons, in our implementation, the numerical solution can naturally maintain its orthogonality when the calculation is carried out to a certain extent. Thus, the orthogonalization operation is avoided in our implementation. Based on this, the time complexity of the whole LOBPCG algorithm can be reduced to $\mathcal{O}(Np)$, which will benefit calculations of large-scale problem significantly.

\begin{remark}
	In this paper, we implement a series of experiments of several molecules. For some molecules with more complex structure, the orthogonalization-free implementation might need to be adjusted. Since there will be some wavefunctions that share the same energy, which makes it is challenging to solve the Kohn-Sham equation accurately.
\end{remark}

\subsection{Parallelization of RR procedure}
Moreover, we accelerate the RR procedure to improve the efficiency. Based on the AFEM framework, the matrices $A$ and $B$ are both sparse. Then the time complexity of the RR procedure is around $\mathcal{O}(Nnz\times p)$, where $nnz$ represents the number of non-zero entries in the sparse matrix ($\ N < Nnz \ll N^2$). To accelerate the RR procedure, we implement it in parallelization, i.e., computation of the matrix-vector multiplications of $A$, $B$, $X$, $W$, and $P$ are implemented parallelly. In this way, the multiplication of $n\times n$ and $n\times p$ matrices is decomposed into the matrix-vector multiplication of the $n\times n$ matrix and a $n\times 1$ vector calculated through $p$ threads. Thus, the time complexity of RR procedure can be reduced to $\mathcal{O}(Nnz)$. 

%\subsection{The efficient implementation implementation of LOBPCG}

	
\section{Numerical experiments}
By avoiding the orthogonalization and implement RR procedure in parallelization, the time complexity of the whole LOBPCG method can be reduced to around $\mathcal{O}(N)$. With the help of the adaptive finite element method, the scale of $N$ can be controlled at a small level. Therefore, our implementation of the LOBPCG in the SCF iteration is quite efficient. Several numerical experiments demonstrate our implementation's feasibility and effectiveness in the following. In this paper, all refer value of ground state energy come from the Computational Chemistry Comparison and Benchmark DataBase (CCCBDB). Based on the AFEPack library, all numerical experiments are implemented on DELL EMC R7525, with 2 AMD EPYC 7713 64-Core Processors, and 900 GB memory.

\subsection{Effectiveness}
%\begin{itemize}
%	\item In the beginning, middle, ending of the SCF iteration, the performance of the LOBPCG method can be shown to illustrate the effectiveness of LOBPCG in the SCF. Shown in Fig \ref{ResinSCF}
%	\item Pick several molecules, and show the convergence of the ground state energy during the SCF iteration. (Methane, Benzene, Azobenzene). Shown in Fig \ref{ResCompare}.
%\end{itemize}
The effectiveness can be demonstrated in the following numerical experiment. Firstly, to show the convergence of the orthogonalization-free implemenetation of LOBPCG method in the SCF iteration, we implement Algorithm \ref{SCFinKSDFT} to compute the ground state energy of the benzene molecule. During the calculation, we record the residual of computed eigenpairs obtained through LOBPCG method in several specified SCF iteration steps, which are plotted in Fig \ref{ResinSCF}.
\begin{figure}[htpb]
	\centering
	\includegraphics[width=0.3\textwidth]{LOBPCGin3SCF}\includegraphics[width=0.3\textwidth]{LOBPCGin16SCF}\includegraphics[width=0.3\textwidth]{LOBPCGin32SCF}
	\caption{The L2-norm residual of eigenvectors in specified SCF iteration steps.}
	\label{ResinSCF}
\end{figure}

The LOBPCG method can always converge to a specified level in different phase of the SCF iteration. Especially, the converged residual become smaller as the SCF iteration progresses. Consequently, the numerical result gets more and more precisely. Based on this, we can demonstrate that the LOBPCG method is effective in SCF iteration. 

Moreover, we implement numerical experiments of methane, benzene. %and azobenzene (still in calculation). 
The numerical result of ground state energy is shown in Fig \ref{ResCompare}.
\begin{figure}[htpb]
	\centering
	\includegraphics[width=0.4\textwidth]{RelativeE4Methane}
	\includegraphics[width=0.4\textwidth]{RelativeE4Ben}
	\caption{The ground state energy of Benzene molecule.}
	\label{ResCompare}
\end{figure}
It can be seen that the relative error (=$\frac{\vert E_h - E_{r}\vert}{E_r}$, where $E_r$ is the refer value) reduces significantly as the mesh grid is adapted. That confirms the effectivity of Algorithm 2 (the SCF framework for solving KS equations with LOBPCG). 

\subsection{Orthogonalization}
%\begin{itemize}
%	\item Fix $p$, showing the linearity corresponding with the mesh number $n$.
%	\item Fix $n$, showing the linearity corresponding with the orbital number $p$.
%\end{itemize}
%More importantly, by avoiding orthogonalization operation, the computational time is saved significantly. The corresponding result of CPU time can be found in Fig \ref{CPUCompare}. It can be seen that as the number of orbitals increases, the CPU time for calculation becomes longer. But the calculations without orthogonalization cost quite less CPU time. The reduction in time consumption becomes more obvious with a large orbital number. As a result, it can be expected that considerable computational time can be saved by avoiding orthogonalization in some more large-scale cases.

%Such two points can be observed in Fig \ref{CPUCompare}.
 Consider the benzene molecule, we implement calculations with and without orthogonalization, respectively. The CPU time per LOBPCG iteration over the number of Dofs is plotted in the left one in Fig \ref{CPUCompare}. The CPU time of the calculation with orthogonalization is always larger than the one without orthogonalization. The difference of CPU time gets larger when more Dofs are added into the mesh grid. 

Then we implement the Algorithm 2 for different molecules on the same mesh grid with and without orhtogonalization, repectively. The CPU time per LOBPCG iteration is recorded in the right one in Fig \ref{CPUCompare}. Calculations without orthogonalization always cost less CPU time. For the molecule with more orbitals, such a reduction in CPU time is more remarkable. Especially, the tendency of the CPU time of calculations with orthogonalization is more affected by the number of orbitals. 

Based on these results, it can be demonstrated that the time complexity of the calculation is reduced from $\mathcal{O}(Np^2)$ to a lower level. 

\begin{remark}
	Ideally, without orthogonality, the time complexity of our implementation can be reduced to $\mathcal{O}(Np)$. Although considerable computational time can be saved by avoiding the orthogonalization, computed eigenpairs may cost much more time during the RR procedure. As a result, the improvement in time consumption may be lower than expected. 
\end{remark}


\begin{figure}[htpb]
	\centering
	\includegraphics[width=0.4\textwidth]{CompareOrthoFixP}\includegraphics[width=0.4\textwidth]{ComparisonOrtho}
	\caption{The CPU time of LOBPCG with and without orthogonalization.}
	\label{CPUCompare}
\end{figure}


\subsubsection{preconditioning}
During numerical experiments, we found that the preconditioner plays an important role in calculations. Without preconditioning, the convergence will be undesirable. We implement numerical experiments of the benzene with Davidson, Elliptic (Poisson) and no preconditioners. The LOBPCG residuals in one SCF iteration step are shown in Figure \ref{ResCompare4pre}. The residuals of the non-preconditioned LOBPCG method tend to oscillate around a specified level, which makes the SCF iterations non-convergent. Although the Davidson and Elliptic precondtioners performs differently, both of them can lead to the convergence of the SCF iteration. 

Thus, to ensure the fast convergence, the Davidson preconditioner and the Poisson preconditioner are used in our implementation. However, both of them have limitations. The Davidson preconditioner is the easy to generate but it only use the diagonal entries of the matrix $A$, which loses many information. As a result, it generally needs more steps to converge as shown in Fig \ref{ResCompare4pre}. As for the Poisson preconditioner, it maintains more information of the original matrix but it usually needs to solve a more complex linear system at Step 8 in Algorithm \ref{SCFinKSDFT}. Consequently, the CPU time per iteration with the Davidson preconditioner is generally shorter. While the LOBPCG method with Poisson preconditioner converges faster. 

To combine the advantages of these two preconditioners, we design a mixing preconditioning strategy in our implementation. In the first few steps of the iteration of the LOBPCG method, we use the Poisson preconditioner. Then the Davidson preconditioner is applied in following iteration steps. This way, the LOBPCG method can converge in fewer steps and with shorter CPU time. To demonstrate that, we implement Algorithm \ref{SCFinKSDFT} to calculate the ground state energy of the benzene molecule on the same mesh grid until the first SCF iteration stops. The CPU time of implementations with different preconditioners is shown in Figure \ref{CPUCompare4pre}. In this implementation, we only use the Poisson preconditioner in the first 5 iteration step of LOBPCG method. 

It can be seen that the mixing strategy can combine the advantages of Davidson and Poisson preconditioners, i.e., it can not only reduce the number of iteration steps but also save computational time. Therefore, it can be confirmed that the mixing strategy can improve efficiency effectively. By now, we mainly implement the mixing strategy based on our numerical experience. A rigorous principle of such a mixing strategy is still in research to fully use its advantages. Furthermore, Poisson preconditioners are pretty helpful in accelerating convergence. But how to reduce the computation time to generate it is an open problem.

\begin{figure}[htpb]
	\centering
	\includegraphics[width=0.5\textwidth]{BenzeneResDiff}
	\caption{CPU time of LOBPCG method in the SCF iteration.}
	\label{ResCompare4pre}
\end{figure}
\begin{figure}[htpb]
	\centering
	\includegraphics[width=0.5\textwidth]{BenzeneCPUDiffPre}
	\caption{Residual of LOBPCG method in the SCF iteration.}
	\label{CPUCompare4pre}
\end{figure}
%\begin{itemize}
%	\item Without preconditioner in LOBPCG, the calculation converges more and more slow.
%	\item With the similar number of iterations, Poisson preconditioner can reach at aroung 1.0e-9. On the contrary, the limit of the Davidson preconditioner is only around 1.0e-8. However, by mixing them up, the accuracy can reach at around 1.0e-9.
%	\item With the similar number of iterations, Poisson preconditioenr costs extremely large CPU time. While the Davidson preconditioner merely spends much less seconds. By mixing them up, the CPU consumption can be controlled to an acceptable level.
%\end{itemize}
%By appropriately using the mixing preconditioning strategy, we can get an accurate enough result in a reasonable CPU time.
%
%In fact, the Hartree potential and Coulomb potential operator might can be added into the Poisson preconditioner, which will be more approximated to the KS Hamiltonian. Due to its similarity to the original matrix, the efficiency of such a preconditioner can be expected \cite{knyazev1998preconditioned}. 
%
%\begin{figure}[htpb]
%	\centering
%	\includegraphics[width=0.5\textwidth]{NewPrecondCompare}
%	\caption{Comparison of the new preconditioner with three preconditioners mentioned above.}
%		\label{CPUCompare4Diff}
%\end{figure}
\subsection{Rayleigh-Ritz procedure}
%\begin{itemize}
%	\item The computational time of RR procedure can be reduced by parallelization as Fig \ref{CPUComparePara}.
%	\item Showing that the time complexity of the Rayleigh-Ritz procedure can be reduced to $1/p$ with the help of parallelization.
%\end{itemize}
To show the acceleration of parallelization, we implement numerical experiments of the benzene molecule in parallelization and in serial firstly. The result is recorded in Fig \ref{CPUComparePara}. It is obvious that the computational time can be saved significantly in the parallelized implementation. When the number of Dofs is huge, the improvement can be more remarkable. Then we compare parallelized and serial implementations of different molecules. The corresponding CPU time is shown in Fig \ref{CPUComparePara2}. The figure illustrates the benefit of parallel implementation, i.e., greatly accelerating the calculation, especially for large-scale molecules. 

\begin{figure}
	\centering
\includegraphics[width=0.5\textwidth]{CompareParallel}
	\caption{The CPU time of LOBPCG implemented parallelly and serially.}
		\label{CPUComparePara}
\end{figure}
\subsection{A numerical example, C60}
%\begin{itemize}
%	\item This is an example of a large molecule to show the capability of our implementation of LOBPCG for SCF iteration.
%\end{itemize}
In this example, we compute the ground state energy of a Buckminsterfullerene (C60) through the Algorithm \ref{SCFinKSDFT}. The computed ground state energy during the iteration and the CPU time for each SCF iteration are both shown in Fig \ref{C60}. It can be observed that the ground state energy get closer to the refer value as the mesh grid is adaptive refined, which demonstrates that the implementation is effective. Furthermore, the relationship between the computational time and the number of Dofs tends to be linear, as shown as the straight line in the figure. Thus, the time complexity of the SCF iteration in Algorithm \ref{SCFinKSDFT} is close to $\mathcal{O}(N)$. which indicates our implementation based on the orthogonalization-free LOBPCG method is quite efficient. 
\begin{figure}[htpb]
	\centering
	\includegraphics[width=0.5\textwidth]{C60}
	\caption{Ground state energy of C60 and corresponding computational time.}
	\label{C60}
\end{figure}

\section{Conclusion}
In this paper, we design an orthogonalization-free implementation of LOBPCG method. Based on this, an adaptive finite element method based framework of SCF iteration for solving the Kohn-Sham equation is established. A series of numerical experiments demonstrate the effectiveness of our implementation. During these experiments, we found a mixing strategy for preconditioning that can further improve the efficiency of the LOBPCG method. Furthermore, considerable improvements in computational time can be observed during calculations. Especially, in the example of C60, the time complexity of SCF iteration can be close to $\mathcal{O}(N)$. With the help of adaptive finite element methods, our implementation can be significantly efficient for calculating the ground state energy. In the future work, we will further improve the efficiency of other parts in the SCF iteration, making the time complexity of the whole SCF iteration tends to $\mathcal{O}(N)$ in theory. 

\section*{Acknowledgement}
The research of G. Hu was partially supported by National Natural Science
Foundation of China (Grant Nos. 11922120 and 11871489), FDCT of Macao SAR
(0082/2020/A2), MYRG of University of Macau (MYRG2020-00265-FST) and
Guangdong-Hong Kong-Macao Joint Laboratory for Data-Driven Fluid Mechanics and
Engineering Applications (2020B1212030001).

\section*{Statements and Declarations}
\subsection*{Data availability statements}
Data sharing not applicable to this article as no datasets were generated or 
analysed during the current study.

\subsection*{Competing interests}
The authors declare that they have no conflict of interest.

\bibliographystyle{plain}
\bibliography{bibliography}
\end{document}
